\documentclass[11pt]{article}
\usepackage[margin=1in]{geometry}
\usepackage{amsmath,amssymb,amsthm,mathtools}
\usepackage{booktabs}
\usepackage{graphicx}
\usepackage{microtype}
\usepackage[hidelinks]{hyperref}

\newtheorem{theorem}{Theorem}
\newtheorem{lemma}[theorem]{Lemma}
\newtheorem{corollary}[theorem]{Corollary}
\newtheorem{remark}[theorem]{Remark}
\newcommand{\R}{\mathbb R}
\newcommand{\vol}{\operatorname{vol}}
\newcommand{\conv}{\operatorname{conv}}

\title{A Half-Quadratic Threshold for Volume Monotonicity\\
of Star-Shaped Minkowski Sums}
\author{Automated Math Discovery Run \texttt{fa\_banach\_001}, lane 4}
\date{9 August 2026}

\begin{document}
\maketitle

\begin{abstract}
For a compact star-shaped set $S\subset\R^d$, Fradelizi, L\'angi, and
Zvavitch proved
\[
 \vol\!\left(\frac{S[k+1]}{k+1}\right)
 \ge \vol\!\left(\frac{S[k]}{k}\right)
 \qquad\text{when }k\ge(d-1)(d-2).
\]
We show that their argument in fact proves the same conclusion, with the same
equality characterization, whenever
\[
       (k-d+2)(k+1)^{d-1}>k^d, \qquad k\ge d-1.
\]
In particular it is enough that
\[
                         k\ge \binom d2-1.
\]
This replaces a threshold asymptotic to $d^2$ by one asymptotic to
$d^2/2$.  The exact scalar condition is slightly stronger still: its first
valid values for $d=4,5,6,7$ are respectively $k=5,8,13,19$, whereas the
published thresholds are $6,12,20,30$.
\end{abstract}

\section{Source problem and scope}

For a set $A\subset\R^d$, write
\[
 A[k]=\underbrace{A+\cdots+A}_{k\text{ summands}}.
\]
The Bobkov--Madiman--Wang conjecture asks whether
$\vol(A[k]/k)$ is nondecreasing for every compact $A$.  The conjecture is
false in sufficiently high dimension, but important structured cases remain
open.  The source paper proves the star-shaped case for $d=2,3$ and, in
arbitrary dimension, for $k\ge(d-1)(d-2)$; its Conjecture~1, definition, and
Theorem~1 appear together on page~2 in Figure~\ref{fig:source}.

\begin{figure}[p]
\centering
\includegraphics[width=0.96\textwidth]{figures/open_problem_crop.png}
\caption{Source page 2: the BMW volume-monotonicity conjecture and the
published star-shaped threshold, arXiv:1910.06146v2.}
\label{fig:source}
\end{figure}

Our result does not settle the remaining low-$k$ star-shaped cases, nor the
general conjecture.  It is a uniform strengthening of the source's positive
range in every dimension $d\ge4$.

\section{Proof intuition}

The source proof decomposes a polytope into simplicial cones and reduces the
problem to adjoining a union $B$ of $d$ radial segments to a set $M$ trapped
between $B[k]$ and $k\conv(B)$.  Its cubical-layer transport estimate gives
the stronger affine inequality
\[
 \vol\!\left(\frac{M+B}{k+1}\right)
 \ge \alpha_{d,k}\vol\!\left(\frac{M}{k}\right)
 +(1-\alpha_{d,k})\vol(\conv B),
 \qquad
 \alpha_{d,k}=\frac{k^d}{(k-d+2)(k+1)^{d-1}}.
\]
Only the last line of that proof replaces the exact requirement
$\alpha_{d,k}<1$ by the convenient sufficient bound
$k\ge(d-1)(d-2)$.  Retaining the exact coefficient immediately enlarges the
range.  A second-order Taylor estimate shows that $\alpha_{d,k}<1$ already
when $k+1\ge\binom d2$.  No geometric part of the argument changes.

\section{The sharpened theorem}

\begin{theorem}[Exact coefficient range]\label{thm:main}
Let $d\ge3$ and $k\ge\max\{2,d-1\}$ be integers satisfying
\begin{equation}\label{eq:delta}
                 (k-d+2)(k+1)^{d-1}>k^d.
\end{equation}
Then every compact star-shaped set $S\subset\R^d$ satisfies
\begin{equation}\label{eq:mono}
 \vol\!\left(\frac{S[k+1]}{k+1}\right)
 \ge \vol\!\left(\frac{S[k]}{k}\right).
\end{equation}
Equality holds if and only if either $\dim S<d$ or
$S[k]/k=\conv S$.
\end{theorem}

\begin{corollary}[Half-quadratic closed form]\label{cor:clean}
The conclusion of Theorem~\ref{thm:main} holds for every $d\ge3$ and
\[
                              k\ge\binom d2-1.
\]
Thus the published condition $(d-1)(d-2)$ can be replaced by
$\binom d2-1$.
\end{corollary}

It is useful to distinguish the exact cutoff for the coefficient
\eqref{eq:delta} from a possibly smaller threshold obtainable by a different
geometric method.  Define
\[
 \kappa_d=\min\{k\ge d-1:(k-d+2)(k+1)^{d-1}>k^d\}.
\]
Exact integer arithmetic gives:
\[
\begin{array}{c|rrrrrrrrrr}
d&3&4&5&6&7&8&9&10&11&12\\ \hline
\kappa_d&2&5&8&13&19&25&33&42&51&62\\
\binom d2-1&2&5&9&14&20&27&35&44&54&65\\
(d-1)(d-2)&2&6&12&20&30&42&56&72&90&110
\end{array}
\]

\section{Local coefficient lemma}

Let $p_1,\ldots,p_d$ be a basis of $\R^d$, let the common endpoint be the
origin, and put
\[
 B=\bigcup_{i=1}^d[0,p_i],\qquad V=\vol(\conv B).
\]

\begin{lemma}[Range-sharpened local lemma]\label{lem:local}
Assume $k\ge\max\{2,d-1\}$, $B[k]\subset M\subset k\conv B$, and $M$ is
compact.  Set
\[
                 \alpha_{d,k}=\frac{k^d}{(k-d+2)(k+1)^{d-1}}.
\]
Then
\begin{equation}\label{eq:affine}
 \vol\!\left(\frac{M+B}{k+1}\right)
 \ge \alpha_{d,k}\vol\!\left(\frac{M}{k}\right)
 +(1-\alpha_{d,k})V.
\end{equation}
If $\alpha_{d,k}<1$, then
\[
 \vol\!\left(\frac{M+B}{k+1}\right)
 \ge \vol\!\left(\frac{M}{k}\right),
\]
with equality only for $M=k\conv B$.  Moreover, if the difference between
the left and right normalized volumes is at most $\delta$, then
\begin{equation}\label{eq:stability}
 \vol(k\conv B)-\vol(M)
 \le \frac{k^d}{1-\alpha_{d,k}}\,\delta.
\end{equation}
\end{lemma}

\begin{proof}
Equation~\eqref{eq:affine} is exactly inequality (10) in the source
paper~\cite{FLZ}, before its final scalar estimate.  We record why the larger
range causes no hidden boundary problem.  The proof partitions $k\conv B$
into unit cubical layers indexed from $t=k-d+1$ to $k-1$.  Every preceding
counting and mass-transport identity is valid as soon as $k-d+1\ge0$; at the
endpoint $k=d-1$, the first layer is simply $t=0$.  The only later use of
$k\ge(d-1)(d-2)$ is the displayed binomial lower bound intended to make
$1-\alpha_{d,k}$ nonnegative.  Thus the derivation of \eqref{eq:affine}
itself is valid for $k\ge d-1$.

Write $u=\vol(M/k)$.  Since $M\subset k\conv B$, one has $u\le V$, and
\eqref{eq:affine} becomes
\[
 \vol\!\left(\frac{M+B}{k+1}\right)
 \ge u+(1-\alpha_{d,k})(V-u).
\]
When $\alpha_{d,k}<1$, this proves monotonicity.  Equality forces $u=V$.
Hence $k\conv B\setminus M$ has measure zero.  The compactness of $M$ and
the fact that a full-dimensional compact convex body is the closure of its
interior imply $M=k\conv B$.  The same displayed inequality gives
$(1-\alpha_{d,k})(V-u)\le\delta$; multiplication by $k^d$ proves
\eqref{eq:stability}.
\end{proof}

\section{Proof of the global theorem}

\begin{proof}[Proof of Theorem~\ref{thm:main}]
Condition~\eqref{eq:delta} is precisely $\alpha_{d,k}<1$, so
Lemma~\ref{lem:local} applies.  Translate a star center of $S$ to the origin.
The case $\dim S<d$ is immediate, so assume $S$ is full-dimensional.

Fix $\varepsilon>0$.  Choose a finite set $A_0\subset S$, enlarged to contain
the origin, with
\[
 \vol(\conv S)-\vol(\conv A_0)\le\varepsilon.
\]
Triangulate the boundary of $\conv A_0$.  Coning its facets to the origin
gives simplices $T_j=\conv B_j$ with mutually disjoint interiors, where each
$B_j$ is the union of the $d$ radial segments from the origin to the vertices
of the corresponding boundary simplex.  Star-shapedness gives $B_j\subset S$.
Put
\[
                         M_j=S[k]\cap kT_j.
\]
Then $B_j[k]\subset M_j\subset kT_j$, so Lemma~\ref{lem:local} yields
\[
 \vol\!\left(\frac{M_j+B_j}{k+1}\right)
 \ge \vol\!\left(\frac{M_j}{k}\right).
\]
The sets on the left lie in $S[k+1]/(k+1)$ and have disjoint interiors.
Since $S[k]/k\subset\conv S$, the portion of $S[k]/k$ outside
$\conv A_0$ has volume at most $\varepsilon$.  Summing over $j$ gives
\[
 \vol\!\left(\frac{S[k]}{k}\right)-\varepsilon
 \le\sum_j\vol\!\left(\frac{M_j}{k}\right)
 \le\sum_j\vol\!\left(\frac{M_j+B_j}{k+1}\right)
 \le\vol\!\left(\frac{S[k+1]}{k+1}\right).
\]
Letting $\varepsilon\downarrow0$ proves \eqref{eq:mono}.

For the equality case, suppose the two endpoint volumes are equal.  Let
$\delta_j$ be the nonnegative difference between the two normalized local
volumes.  The preceding chain gives $\sum_j\delta_j\le\varepsilon$.
By \eqref{eq:stability}, with
$C_{d,k}=k^d/(1-\alpha_{d,k})$,
\[
 \vol(k\conv A_0)-\vol(S[k]\cap k\conv A_0)
 \le C_{d,k}\varepsilon.
\]
The part of $k\conv S$ outside $k\conv A_0$ has volume at most
$k^d\varepsilon$.  Consequently
\[
       \vol(k\conv S)-\vol(S[k])
       \le (k^d+C_{d,k})\varepsilon.
\]
Letting $\varepsilon\downarrow0$ and using compactness as in the local lemma
gives $S[k]=k\conv S$.  Conversely, if this set equality holds, the left
normalized volume is already $\vol(\conv S)$; monotonicity and the containment
$S[k+1]/(k+1)\subset\conv S$ force equality.  This proves the stated
characterization.
\end{proof}

\section{The half-quadratic scalar estimate}

\begin{proof}[Proof of Corollary~\ref{cor:clean}]
Put $x=k+1$ and $t=1/x$.  Condition~\eqref{eq:delta} is equivalent to
\begin{equation}\label{eq:h}
                  (1-t)^d<1-(d-1)t.
\end{equation}
For $d\ge3$ and $0<t<1$, Taylor's theorem at zero gives the strict upper
bound
\[
                 (1-t)^d<1-dt+\binom d2t^2,
\]
because the third derivative of $(1-t)^d$ is negative.  If
$x\ge\binom d2$, then $\binom d2t\le1$, and hence
\[
 1-dt+\binom d2t^2\le1-dt+t=1-(d-1)t.
\]
This proves \eqref{eq:h}, hence \eqref{eq:delta}.  Since
$k\ge\binom d2-1$ also implies $k\ge d-1$, Theorem~\ref{thm:main} applies.
\end{proof}

For completeness, the exact coefficient cutoff is unambiguous.  The function
\[
                h_d(t)=1-(d-1)t-(1-t)^d
\]
has $h_d(0)=0$, $h_d'(0)=1$, and
$h_d''(t)=-d(d-1)(1-t)^{d-2}<0$ on $(0,1)$.  Thus it has at most one
positive zero, and the positivity test can change only once as $k$ increases
and $t=1/(k+1)$ decreases.  The table above therefore records exact cutoffs
for this coefficient argument, not a collection of isolated successes.

\section{Verification, novelty, and limitations}

The reusable checker in \texttt{code/verify\_thresholds.py} uses exact integer
arithmetic.  It verifies the displayed values of $\kappa_d$, confirms strict
positivity at the closed-form threshold through $d=1000$, and checks the
equivalence between the polynomial and normalized forms of the scalar
condition through $d=200$.  These computations are sanity checks; the proof
of the uniform statement is the Taylor argument above.

The main human-review point is narrow: confirm that the cubical-layer
derivation of source inequality (10) uses only $k\ge d-1$ before the final
coefficient estimate.  The source text has been audited term by term, including
the endpoint layer $t=0$.

For novelty, the run's solution, attempt, and proof-gap indexes were searched
for arXiv:1910.06146 and the star-shaped threshold.  Focused arXiv searches on
9 August 2026 used the exact paper title, the scalar factor
\texttt{(k-d+2)(k+1)}, and combinations of ``star-shaped'', ``Minkowski
sum'', ``volume monotonicity'', and ``threshold''.  They found the source
paper~\cite{FLZ}, the earlier general counterexample paper~\cite{FMMZ}, and
broader subset-sum work, but no later paper stating this refinement.  This is
a bounded novelty check, not a claim of exhaustive bibliographic priority.

The packet is a substantial partial result.  It does not resolve the remaining
star-shaped cases $k<\kappa_d$; for example, dimensions four and five still
have the unresolved ranges $k=2,3,4$ and $k=2,3,4,5,6,7$, respectively.

\begin{thebibliography}{9}
\bibitem{FLZ}
M. Fradelizi, Z. L\'angi, and A. Zvavitch,
\emph{Volume of the Minkowski sums of star-shaped sets},
arXiv:1910.06146v2 (2021),
\url{https://arxiv.org/abs/1910.06146}.

\bibitem{FMMZ}
M. Fradelizi, M. Madiman, A. Marsiglietti, and A. Zvavitch,
\emph{Do Minkowski averages get progressively more convex?},
arXiv:1512.03718 (2015),
\url{https://arxiv.org/abs/1512.03718}.
\end{thebibliography}

\end{document}
